% Figure: pendulum period as a function of amplitude. Compares the
% small-angle prediction, the first-correction prediction, and the
% exact (elliptic-integral) result.
\begin{figure}[htbp]
\centering
\begin{tikzpicture}[
  axis/.style={draw=engblack, -{Stealth}, thick},
  small/.style={draw=engblue, very thick, dashed},
  first/.style={draw=engamber, very thick, dash dot},
  exact/.style={draw=engred, very thick},
  guide/.style={draw=enggray, dashed, thin},
  label/.style={font=\small},
]
% Axes: x amplitude in degrees (0 to 120, plotted x in 0..10), y T/T0 (1.0 to 1.4, plotted 0..5)
\draw[axis] (0,0) -- (10.5,0) node[anchor=west, label] {$\theta_{0}$ [deg]};
\draw[axis] (0,0) -- (0,5.5)  node[anchor=south, label] {$T/T_{0}$};

% x ticks (12 deg per unit)
\foreach \x/\m in {0/0, 2/24, 4/48, 6/72, 8/96, 10/120} {
  \draw (\x,-0.08) -- (\x,0.08);
  \node[font=\scriptsize, below] at (\x,-0.1) {\m};
}
% y ticks (0.08 per unit; 1.0 at y=0; 1.40 at y=5)
\foreach \y/\h in {0/1.00, 1.25/1.10, 2.5/1.20, 3.75/1.30, 5/1.40} {
  \draw (-0.08,\y) -- (0.08,\y);
  \node[font=\scriptsize, left] at (-0.1,\y) {\h};
}

% Reference line at T/T0 = 1.0
\draw[guide] (0, 0) -- (10, 0);
% Reference line at T/T0 = 1.05 (y = 0.625)
\draw[guide] (0, 0.625) -- (10, 0.625);
\node[engblack, font=\scriptsize, anchor=west] at (10.05, 0.625) {$+5\,\%$};

% Small-angle prediction: T/T0 = 1 always.
\draw[small] (0, 0) -- (10, 0);
\node[engblue, font=\scriptsize, anchor=west] at (10.05, 0) {small-angle};

% First correction: T/T0 = 1 + theta^2/16, theta in radians
% x = theta_deg / 12, theta_deg = 12*x, theta_rad = 12*x*pi/180 = x*pi/15
% Correction = (x*pi/15)^2 / 16; T/T0 - 1 plotted as 12.5 * correction
\draw[first] plot[domain=0:10, samples=120]
  ({\x}, {12.5 * ((\x * 3.14159/15)^2 / 16)});
\node[engamber, font=\scriptsize, anchor=west] at (8.0, 1.7) {first correction $1 + \theta_{0}^{2}/16$};

% Exact (elliptic-integral) values: T/T0 = (2/pi) K(sin(theta/2))
% Sampled at theta_deg = 0, 10, 20, ..., 120
% K(k) values: theta/2 in degrees, k = sin(theta/2 in rad)
% Reference values (Abramowitz & Stegun):
% theta=0:    T/T0 = 1.00000
% theta=10:   T/T0 = 1.00191
% theta=20:   T/T0 = 1.00767
% theta=30:   T/T0 = 1.01741
% theta=40:   T/T0 = 1.03135
% theta=50:   T/T0 = 1.04978
% theta=60:   T/T0 = 1.07314
% theta=70:   T/T0 = 1.10200
% theta=80:   T/T0 = 1.13715
% theta=90:   T/T0 = 1.17984
% theta=100:  T/T0 = 1.23198
% theta=110:  T/T0 = 1.29661
% theta=120:  T/T0 = 1.37864
% Translate: x = theta_deg/12, y = 12.5 * (T/T0 - 1)
\draw[exact] plot coordinates {
  (0,      0.000)
  (0.833,  0.024)
  (1.667,  0.096)
  (2.500,  0.218)
  (3.333,  0.392)
  (4.167,  0.622)
  (5.000,  0.914)
  (5.833,  1.275)
  (6.667,  1.714)
  (7.500,  2.248)
  (8.333,  2.900)
  (9.167,  3.708)
  (10.000, 4.733)
};
\node[engred, font=\small, anchor=south west] at (5.5, 3.4) {exact (elliptic)};

% Mark amplitude at which exact crosses 5% (~67 deg). y=0.625 at x ~ 5.58
\draw[guide] (5.58, 0) -- (5.58, 0.625);
\node[font=\scriptsize, enggray, anchor=south west] at (5.5, 0.0) {$\sim 67^{\circ}$};

\end{tikzpicture}
\caption[Pendulum period as a function of amplitude]{The pendulum
period $T(\theta_{0})$ normalised by the small-angle period
$T_{0}=2\pi\sqrt{\ell/g}$, plotted against the amplitude
$\theta_{0}$ in degrees. The small-angle prediction (blue dashed)
is flat at $T/T_{0}=1$. The first amplitude correction
$T/T_{0}\approx 1+\theta_{0}^{2}/16$ (amber dash-dot) tracks the
exact value to within $1\,\%$ up to about $\theta_{0}=70^{\circ}$
and diverges thereafter. The exact value, obtained from the
complete elliptic integral of the first kind $K$, follows the red
curve. The amplitude at which the exact period exceeds the
small-angle period by $5\,\%$ is approximately $67^{\circ}$;
beyond $\theta_{0}\approx 90^{\circ}$ the small-angle formula
under-predicts the period by more than $18\,\%$. The chapter's
project quantifies the crossings; the figure provides the
reference picture against which numerical extractions can be
sanity-checked. Exact values from Abramowitz \& Stegun, table 17.1.}
\label{fig:vol02:ch07:pendulum-period}
\end{figure}
